%!TEX root = ../dissertation.tex
% !TeX spellcheck = en_GB 

\chapter{The Algorithm}\label{chapt:2}
%This chapter is divided in four sections: Section~\ref{sect:2_1}, Section~\ref{sect:2_2}, Section~\ref{sect:2_3}.\\

In this chapter the proposed algorithm is described in section \ref{sect:2_4} while the previous sections explain the various pieces that compose it.

%----------------------------------------------------

\section{The Gradient Estimator}

\subsection{Gaussian Smoothing}

The standard method used to estimate the gradient $\nabla f(x)$, which was formalised by \cite{nesterov2017random} and used in the further literature on the subject, is the Gaussian Smoothing. Let's introduce it.

\begin{defn}[Gaussian Smoothing]\label{GaussianSmoothingDefn}

	Let $\sigma > 0$ the \textit{smoothing parameter}, then we denote by $f_\sigma (x)$ the Gaussian Smoothing of $f$ with radius $\sigma$, i.e.
	
	\begin{equation}
	    f_{\sigma}(x)=\frac{1}{\pi^{d / 2}} \int_{\mathbb{R}^{d}} f(x+\sigma \xi) e^{-\|\xi\|_{2}^{2}} \mathrm{~d} \xi=\mathbb{E}_{\xi \sim \mathcal{N}\left(0, \mathbb{I}_{i}\right)}[f(x+\sigma \xi)]
	\end{equation}

\end{defn}

Note that $f_\sigma$ maintains features of $f$ (convexity and Lipschitz constant) while being always differentiable, even when $f$ is not. This is relevant because, given that $\norm{f - f_\sigma}$ can be bound by the Lipschitz constant, \ref{problem} can be substituted with

\begin{equation}\label{smoothed problem}
    \min _{x \in \mathbb{R}^{d}} f_\sigma (x)  
\end{equation}

Now, the gradient of the smoothed $f$ is:

\begin{equation}
    \nabla f_{\sigma}(x)=\frac{2}{\sigma \pi^{d / 2}} \int_{\mathbb{R}^{d}} \xi f(x+\sigma \xi) e^{-\|\xi\|_{2}^{2}} \mathrm{~d} \xi=\frac{2}{\sigma} \mathbb{E}_{\xi \sim \mathcal{N}\left(0, \mathbb{I}_{d}\right)}[\xi f(x+\sigma \xi)]
\end{equation}

The estimator usually exploited to estimate this gradient is obtained via Monte Carlo sampling, with a central finite difference approach, which is used to reduce the variance of the estimation; namely, given $M \in \mathbb{N}_+$ the number of directions and $\left\{\varepsilon_{j}\right\}_{j=1}^{M} \stackrel{\mathrm{iid}}{\sim} \mathcal{N}(0,1)$ the directions, then we have the following:

\begin{defn}[Gradient of Gaussian Smoothed function]\label{GradientGaussianSmoothingDefn}

	\begin{equation}
	    \nabla f_{\sigma}(x) \approx \frac{1}{\sigma M} \sum_{j=1}^{M} \xi{j}\left(f\left(x+\sigma \xi{j}\right)-f\left(x-\sigma \xi{j}\right)\right)
	\end{equation}

\end{defn}

\subsection{Directional Gaussian smoothing via Gauss-Hermite quadrature}

The first idea of the algorithm here presented, is, instead of doing a single $d-$dimensional Gaussian Smoothing, to perform $d$ $1-$dimensional Gaussian Smoothing, which can be approximated via Gauss-Hermite quadrature. 

To achieve this, a function $G:\mathbb{R} \rightarrow \mathbb{R}$ is defined as a cross section of $F$ along $\xi$ as

\begin{equation}
 G(y \mid x, \xi)=F(x+y \xi), \quad y \in \mathbb{R} 
\end{equation}

It is further defined the Gaussian smoothing of $G(y)$, denoted by $G_\sigma (y)$, by

\begin{equation}
 G_{\sigma}(y \mid x, \xi):=\frac{1}{\sqrt{2 \pi}} \int_{\mathbb{R}} G(y+\sigma v \mid x, \xi) \mathrm{c}^{-\frac{r^{2}}{2}} d v=\mathbb{E}_{v \sim \mathcal{N}(0,1)}[G(y+\sigma v \mid x, \xi)] 
\end{equation}

which is the Guassian smoothing of $F(x)$ along $\xi$ near $x$. Using a one-dimensional expectation, the derivative of $G(y \mid x, \xi)$ at $y=0$, is 

\begin{equation}\label{DerivativeSmoothingSection}
 \mathcal{D}\left[G_{\sigma}(0 \mid x, \xi)\right]=\frac{1}{\sigma} \mathbb{E}_{v \sim \mathcal{N}(0,1)}[G(\sigma v \mid x, \xi) v] 
\end{equation}

Now let $\Xi = \left(\xi_1,\dots,\xi_d\right)$ an orthonormal basis, than the following vector, referenced at as DGS gradient, can be defined:

\begin{defn}[DGS Gradient]\label{DGSGradient}

	\begin{equation}
		\nabla_{\sigma, \xi}[F](x):=\left[\mathcal{D}\left[G_{\sigma}\left(0 \mid x, \xi_{1}\right)\right], \cdots, \mathcal{D}\left[G_{\sigma}\left(0 \mid x, \xi_{d}\right)\right]\right] \xi 
	\end{equation}

\end{defn}

At this point, the key is that each of the components of the DGS gradient requires a one-dimensional integral, which is approximated with high accuracy with the Gauss-Hermite rule. Hence we obtain the following estimator for $\mathcal{D}\left[G_{\sigma}(0 \mid x, \xi)\right]$

\begin{equation}\label{GHApproximationSmothedSection}
	\tilde{\mathcal{D}}^{M}\left[G_{\sigma}(0 \mid x, \xi)\right]=\frac{1}{\sqrt{\pi} \sigma} \sum_{m=1}^{M} w_{m} F\left(x+\sqrt{2} \sigma v_{m} \xi\right) \sqrt{2} v_{m} 
\end{equation}

Thus we define the DGS Gradient Estimator as 

\begin{defn}[DGS Gradient Estimator]\label{DGSGradientEstimator}

	\begin{equation}\label{DGSGradientEstimatorEq}
		\widetilde{\nabla}_{\sigma, \Xi}^{M}[F](x)=\left[\widetilde{\mathcal{D}}^{M}\left[G_{\sigma}\left(0 \mid x, \xi_{1}\right)\right], \cdots, \widetilde{\mathcal{D}}^{M}\left[G_{\sigma}\left(0 \mid x, \xi_{d}\right)\right]\right] \Xi 
	\end{equation}

\end{defn}

%----------------------------------------------------

\section{Historical Gradients}\label{sect:2_2}

The second core idea of the algorithm is the use of gradients computed in previous iterations in order to identify a better set of directions along which estimate the gradient at the current iteration. To do so, $k\ll d$ gradients are collected from the beginning, and then always the most recent $k$ gradients in a buffer are kept. Then at each iteration $t \geq k$, given $\mathbf{G}_t\in \mathbb{R}^{n\text{x} k}$, which is a matrix made with the last $k$ gradients, the algorithm generates the subspace $\mathcal{L}_\mathbf{G}$, using the vectors in the buffer. This subspace captures the structure of the optimization problem and, in the reinforcement learning optic, it captures the directions related to exploitation.
Then the algorithm samples the search directions from:

\begin{equation}\label{historical sampling}
	\left\{\begin{array}{ll}\xi \sim \mathcal{N}\left(0, \mathrm{G}^{\perp} \mathrm{G}\right) & \text { with probability } \alpha \\ \xi \sim \mathcal{N}(0, \mathrm{I}) & \text { with probability } 1-\alpha\end{array}\right\}. 
\end{equation}

$\alpha\in\left( 0,1 \right)$ represents the trade-off between exploitation, choosing a direction from $L_\mathbf{G}$, and exploration, from the whole space. The algorithm adapts $\alpha$ at each iteration such that the amount of exploitation and exploration varies depending on how the algorithm is performing. To do that the following are computed:

\begin{equation}
	\widehat{r}_{\mathbf{G}}=\frac{1}{M} \sum_{i=1}^{M} \min_{m=1,\dots,m_i} \left\{ f\left( x+\sigma p_m \xi_i \right) \right\} 
\end{equation}

\begin{equation}
	\widehat{r}_{\mathbf{G}}^{\perp}=\frac{1}{d-M} \sum_{i=M+1}^{d} \min_{m=1,\dots,m_i} \left\{ f\left( x+\sigma p_m \xi_i \right) \right\} 
\end{equation}

Then, at iteration $t$, $\alpha$ is chosen as:

\begin{equation}\label{adapt alpha}
	\alpha_{t}=\left\{\begin{array}{ll}\min \left\{\delta \alpha_{t-1}, \kappa_{1}\right\} \text { if } \widehat{r}_{\mathbf{G}} \text { does not exist or } \widehat{r}_{\mathbf{G}}<\widehat{r}_{\mathbf{G}}^{\perp} \\ \max \left\{\frac{1}{\delta} \alpha_{t-1}, \kappa_{2}\right\} \text { if } \widehat{r}_{\mathbf{G}}^{\perp} \text { does not exist or } \widehat{r}_{\mathbf{G}} \geq \widehat{r}_{\mathbf{G}}^{\perp}\end{array}\right. 
\end{equation}

with $\delta > 1$ a scaling factor, $k_1$ and $k_2$ are upper and lower bounds of $\alpha$. Intuitively, if $\widehat{r}_{\mathrm{G}}<\widehat{r}_{\mathrm{G}}^{\perp} $ or $\widehat{r}_{\mathrm{G}}$ does not exist, than the algorithm increases $\alpha$, otherwise it is decreased.

%----------------------------------------------------

\section{Adaptation of the Parameters}\label{sect:2_3} 

Lastly, instead of keeping for the whole execution of the algorithm some fixed values of the smoothing parameter $\sigma$ and the learning rate $\lambda$, it modifies them at each iteration, in order to estimate their best values.

\subsection{Learning rate}

To update the learning rate, at each iteration the directional local Lipschitz constants $L_j$ are estimated as:

\begin{equation}\label{Lipschitz constants}
	L_{j}=\max _{\{i,k\}\in I}\left|\frac{f\left(x+\sigma p_{i} \xi_{j}\right)-f\left(x+\sigma p_{k} \xi_{j}\right)}{\sigma\left(p_{i}-p_{k}\right)}\right| 
\end{equation}

where $I=\set[\big]{ \{ i,k \} \in \{1,\dots,m \}^2 \given \mid i-\floor{\frac{m}{2}} -1\mid \neq \mid k-\floor{\frac{m}{2}} -1\mid }$. Then the learning rate is computed as:

\begin{equation}\label{learning rate}
     L_{\nabla} \leftarrow\left(1-\gamma_{L}\right) L_{\mathbf{G}}+\gamma_{L} L_{\nabla} \quad \text{and} \quad \lambda=\sigma / L_{\nabla}
\end{equation}

where $L_{\mathbf{G}}$ is the max of the estimated Lipschitz constants associated to the directions sampled from $\mathcal{L}_{\mathbf{G}}$. Note that until iteration $T_w$, $L_{\mathbf{G}}$ is the max of all the local Lipschitz constants.

\subsection{Smoothing parameter}
 
At each iteration, to adapt the smoothing parameter to the local geometry, $\forall j=1,\dots,d$, $\widetilde{\mathcal{D}}^{M}\left[G_{\sigma}\left(0 \mid x, \xi_{j}\right)\right]$ is compared to the corresponding local Lipschitz constant $L_j$. If the max of all the ratios of the module of the two quantities is sufficiently large, $\sigma$ is increase, else it is increased.

Furthermore, the algorithm as a given number of reset of the smoothing parameter $\sigma$ and directions to initial conditions. This is implemented in order to help the algorithm escaping local minima that would otherwise make it stop. This happens whenever the smoothing parameter becomes sufficiently small, $\sigma < \rho \sigma_0$.

%----------------------------------------------------

\section{The Algorithm}\label{sect:2_4} 

\begin{algorithm}[!ht]
\caption{ASHGF algorithm}\label{ASHGF}

\hspace*{\algorithmicindent} \textbf{Input:} function $f$, state $x_0$, smoothing parameter $\sigma_0$, termination tolerance $\varepsilon$, hyper-parameters $\alpha, k$, total iterations \textit{maxiter}, warmup iteration $T_w = k$\\
\hspace*{\algorithmicindent} \textbf{Output:} point of minimum $\Tilde{x}$
\begin{algorithmic}[1]

\State Set $\Tilde{x} = x_0$, $\sigma = \sigma_0$ and let $\Xi$ be a random orthonormal basis
\State Initialize an archive $\mathcal{G}$ of maximum capacity $k$
\For{$i=0$ to \textit{maxiter}}
    \For{$j=1$ to $d$}
        \State compute $ \widetilde{\mathcal{D}}^{M}\left[G_{\sigma}\left(0 \mid x, \xi_{j}\right)\right] $ by \ref{GHApproximationSmothedSection} and estimate local Lipschitz constants $L_j$ by \ref{Lipschitz constants}
    \EndFor
    \State Average Lipschitz constant $L_\nabla$ and compute learning rate $\lambda$ by \ref{learning rate}
    \State Assemble the gradient $ \widetilde{\nabla}_{\sigma, \Xi}^{M}[F](x_i) $ by \ref{DGSGradientEstimatorEq} and update $ x_{i+1}=x_{i}-\lambda \widetilde{\nabla}_{\sigma, \Xi}^{M}[F](x_i) $
    \State Add the gradient estimate to $\mathbf{G}$
    \If{$f\left( x_{i+1}\right) < f(\Tilde{x})$}
        \State Update $\Tilde{x} = x_{i+1}$
    \EndIf
    \If{$\norm{x_{i+1} - x_i}_2 <\varepsilon$}
        \State \textbf{break}
    \Else
        \If{$i < T_w$}
            \State Set $historical = False$
        \Else
            \If{$i >= T_w + 1$}
                \State Adaptively adjust $\alpha$ by \ref{adapt alpha}
            \EndIf
            \State Set $historical = True$
        \EndIf
        \State Update smoothness parameter $\sigma$ and the search directions $\Sigma$ by Algorithm \ref{ASHGFSubroutine}
    \EndIf
\EndFor

\end{algorithmic}
\end{algorithm}

The algorithm, shown in \ref{ASHGF}, presents itself as an alteration of the random search defined by \cite{nesterov2017random} and the Evolutionary Strategy by \cite{Salimans2017}. As said in the previous sections, the three main differences, with respect to the algorithms from the literature, are:
\begin{itemize}
    \item the peculiar choice of the directions along which estimate the the gradient, influenced by past gradients;
    \item the adaptivity of the learning rate and of the smoothing parameter;
    \item the use of Gauss-Hermite quadrature instead of finite differences.
\end{itemize}

Furthermore, it is to be noted that the algorithm does not use any kind of technique which stems from the zeroth-order optimization literature. In fact it doesn't exploit variance reduction, momentum, the sign of gradient, estimated hessian or other approaches. It is, in this regard, a simple random search.

\begin{algorithm}[!ht]
\caption{Parameter update}\label{ASHGFSubroutine}

\hspace*{\algorithmicindent} \textbf{Input:} smoothing parameter $\sigma$, gradient $ \widetilde{\nabla}_{\sigma, \Xi}^{M}[F](x_i)$, local Lipschitz constants $L_1,\dots,L_d$, $historical$, archive of gradients $\mathcal{G}$\\
\hspace*{\algorithmicindent} \textbf{Data:} number of resets $r$ and reset factor $\rho$, decay rate $\gamma_\sigma$, threshold parameters $A,B$ and their change rates $A_+, A_-, B_+, B_-$\\
\hspace*{\algorithmicindent} \textbf{Output:} smoothing parameter $\sigma$ and directions $\Sigma$
\begin{algorithmic}[1]

\If{$r>0$ \textbf{and} $\sigma < \rho \sigma_0$}
    \State Assign $\Xi$ to be a random orthonormal basis and set $\sigma = \sigma_0$
    \State Set $A,B$ to their initial values and set $r = r-1$
\Else
    \If{$historical$}
        \State Obtain gradient matrix $\mathbf{G}\in\mathbb{R}^{d\text{x}k}$ from $\mathcal{G}$
        \State Generate $d$ directions as in \ref{historical sampling}
        \State Ortho-normalize the directions $\mathcal{D}$
        \While{$\mathcal{D}$ does not generate $\mathbb{R}^d$}
            \State Complement $\mathcal{D}$ with random vectors from $\mathcal{N}\left(0, \mathbf{I}\right)$ and orthonormalize
        \EndWhile
    \Else
        \State Set as $\mathcal{D}$ a random orthonormal basis
    \EndIf
    \State Update $\Xi = \mathcal{D}$
    \If{\( \max _{1 \leq j \leq d}\left| \widetilde{\mathcal{D}}^{M}\left[G_{\sigma}\left(0 \mid x_i, \xi_{j}\right)\right] / L_{j}\right|<A \)}
        \State Decrease smoothing $\sigma = \sigma * \gamma_{\sigma}$ and lower threshold $A=A*A_-$
    \ElsIf{\( \max _{1 \leq j \leq d}\left| \widetilde{\mathcal{D}}^{M}\left[G_{\sigma}\left(0 \mid 
    	, \xi_{j}\right)\right] / L_{j}\right|> B \)}
        \State Increase smoothing $\sigma = \sigma / \gamma_{\sigma}$ and upper threshold $B=B*B_+$
    \Else
        \State Increase lower threshold $A=A*A_+$ and decrease upper threshold $B=B*B_-$
    \EndIf
\EndIf

\end{algorithmic}
\end{algorithm}

\subsection{Implementation Details}

Below the table with the parameters used can be found:

\begin{table}[!ht]
%\centering
\hspace{-13mm}
\begin{tabular}{|l|llllllllllllll|}
\cline{1-15}
 parameter & $m$ & $A$   & $B$   & $A_-$ & $A_+$ & $B_-$ & $B_+$ & $\gamma_L$ & $\gamma_\sigma$ & $r$  & $\rho$ & $k$ & $maxiter$ & $\varepsilon$\\ \cline{1-15}
 value & $5$ & $0.1$ & $0.9$ & $0.95$ & $1.02$ & $0.98$ & $1.01$ & $0.9$ & $0.9$ & $10$ & $0.01$ & $50$ & $10000$ & $10^{-8}$            \\ \cline{1-15}
\end{tabular}
\caption[Table of the parameters.]{Table of the parameters.}
\end{table}

The strength of this algorithm lies also in the fact that varying these parameters, keep almost the same performance on average, with some fluctuations depending on the function that is to be optimised.

\section{Theoretical Results}

\subsection{Background notions and results}

In this section the needed theoretical background needed in order to prove convergence results for the algorithm is presented.

\begin{defn}[Lipschitz Continuos Gradient]\label{LipschitzContinousGradient}
Let $F\in\mathcal{C}^1 (\mathbb{R})$. $F$ has Lipschitz Continuos Gradient if there exists $L>0$ such that $\norm{\nabla F(x+\xi) - \nabla F(x)} \leq L \norm{\xi}$, $\forall x, \xi \in \mathbb{R}^d$.
\end{defn}

\begin{defn}[Strong Convexity]
Let $F\in\mathcal{C}^1 (\mathbb{R}^d)$. $F$ is a strongly convex function if there exists a positive number $\tau$ such that $\forall x, \xi \in \mathbb{R}^{d}, F(x+\xi) \geq F(x)+\langle\nabla F(x), \xi\rangle+\frac{\tau}{2}\|\xi\|^{2}$. $\tau$ is called the convexity parameter of $F$.
\end{defn}

\begin{thm}[Gauss-Hermite quadrature error]\label{GHError}

Let $\widetilde{\mathcal{D}}^{M}$ and $\mathcal{D}$ as in \ref{GHApproximationSmothedSection} and \ref{DerivativeSmoothingSection}, respectively. Then the error of the one-dimensional GH quadrature is

\begin{equation}\label{BoundErrorGHEq}
\left|\left(\widetilde{\mathcal{D}}^{M}-\mathcal{D}\right)\left[G_{\sigma}\right]\right| \leq C \frac{M ! \sqrt{\pi}}{2^{M}(2 M) !} \sigma^{2 M-1},   
\end{equation}

with $C>0$ a constant independent of $M$ and $\sigma$.

\end{thm}

\begin{prop}[Lemma 3 - Random gradient-free minimization of convex functions]\label{PrimoLemmaNesterov}
If $ F $ ha Lipschitz Continuos Gradient with constant $ L $, then

\begin{equation}
\left\|\nabla F_{\sigma}(x)-\nabla F(x)\right\| \leq \frac{\sigma}{2} L(d+3)^{3 / 2}
\end{equation}

For $ F \in C^{2}(\mathbb{R}^d) $ such that $\norm{\nabla^2 F(x+\xi) - \nabla^2 F(x)} \leq L \norm{\xi}$, $\forall x, \xi \in \mathbb{R}^d$ with constant $ L $, we can guarantee that

\begin{equation}
\left\|\nabla F_{\sigma}(x)-\nabla F(x)\right\| \leq \frac{\sigma^{2}}{6} L(d+4)^{2}
\end{equation}

\end{prop}

\begin{prop}[Theorem 2.1.5 - Introductory Lectures on Convex Optimization]\label{TeoremaNesterov}
All conditions below, holding for all $ x, y \in R^{n} $ and $ \alpha $ from $ [0,1] $, are equivalent to to saying that $F$ has Lipschitz Continous Gradient
\begin{equation}
0 \leq F(y)-F(x)-\left\langle F^{\prime}(x), y-x\right\rangle \leq \frac{L}{2}\norm{x-y}^{2},
\end{equation}

\begin{equation}
F(x)+\left\langle F^{\prime}(x), y-x\right\rangle+\frac{1}{2 L}\norm{F^{\prime}(x)-F^{\prime}(y)^{2}} \leq F(y),    
\end{equation}

\begin{equation}
\frac{1}{L}\norm{F^{\prime}(x)-F^{\prime}(y)}^{2} \leq\left\langle F^{\prime}(x)-F^{\prime}(y), x-y\right\rangle,    
\end{equation}

\begin{equation}
\left\langle F^{\prime}(x)-F^{\prime}(y), x-y\right\rangle \leq L\norm{x-y}^{2},  
\end{equation}

\begin{equation}
\alpha F(x)+(1-\alpha) F(y) \geq F(\alpha x+(1-\alpha) y) + \frac{\alpha(1-\alpha)}{2 L}\norm{F^{\prime}(x)-F^{\prime}(y)}^{2},   
\end{equation}

\begin{equation}
\alpha F(x)+(1-\alpha) F(y) \leq F(\alpha x+(1-\alpha) y) +\alpha(1-\alpha) \frac{L}{2}\norm{x-y}^{2}.   
\end{equation}

\end{prop}

\subsection{Convergence of the algorithm}

A first result is a bound on the difference between $\nabla F$ and the DGS estimator.

\begin{prop}\label{DGSEstimatorBound}

Let $\Xi = \left(\xi_1,\dots,\xi_d\right)$ be a orthonormal basis of $\mathbb{R}^d$ and let $F$ have a Lipschitz Continous Gradient with constant $L$. Then

\begin{equation}\label{BoundErrorEstimatorEq}
 \left\|\widetilde{\nabla}_{\sigma, \Xi}^{M}[F](x)-\nabla F(x)\right\|^{2} \leq \frac{2 C^{2} \pi d(M !)^{2}}{4^{M}((2 M) !)^{2}} \sigma^{4 M-2}+32 d L^{2} \sigma^{2} 
\end{equation}

\end{prop}

\begin{proof}

First, adapting \ref{PrimoLemmaNesterov} to $1-$dimensional Gaussian smoothing, for any $ \xi $ being a unit vector in $ \mathbb{R}^{d} $, there holds

\begin{equation}\label{BoundOnErrorGradientSmoothingSectionEq}
\left|\mathcal{D}\left[G_{\sigma}(0 \mid x, \xi)\right]-\nabla_{\xi} F(x)\right| \leq 4 \sigma L,
\end{equation}

where $d=1$, being a $1-$dimensional Gaussian smoothing.

From \ref{BoundErrorGHEq} and \ref{BoundOnErrorGradientSmoothingSectionEq}, we have

\begin{equation}\label{BoundOnErrorEstimationGradientSmoothingSectionEq}
 \begin{aligned}\left|\widetilde{\mathcal{D}}^{M}\left[G_{\sigma}\left(0 \mid x, \xi_{i}\right)\right]-\nabla_{\xi_{i}} F(x)\right|^{2} & \leq 2\left|\left(\widetilde{\mathcal{P}}^{M}-\mathcal{D}\right)\left[G_{\sigma}\right]\right|^{2}+2\left|\mathcal{D}\left[G_{\sigma}\right]-\nabla_{\xi_{i}} F(x)\right|^{2} \\ & \leq \frac{2 C^{2}(M !)^{2} \pi}{4^{M}((2 M) !)^{2}} \sigma^{4 M-2}+32 \sigma^{2} L^{2}, \forall i \in\{1, \ldots, d\} \end{aligned} 
\end{equation}

Summing \ref{BoundOnErrorEstimationGradientSmoothingSectionEq} from $ i=1 $ to $ d $, given that the bound doesn't depend on $i$, gives \ref{BoundErrorEstimatorEq}.

\end{proof}

The next result aims to give a bound on the difference in the value of the function at an iteration $t$ with respect to its global minima.

\begin{thm}[Bound on the descent]\label{DescentBoundThm}

Let $ F $ be a strongly convex function with Lipschitz Continuos Gradient, $L$ being its Lipschitz constant, $x^*$ the global minimizer of $ F $ and the sequence $ \left\{x_{t}\right\}_{t \geq 0} $ be generated by Algorithm \ref{ASHGF} with $ \lambda=DATROVARE $. Then, for any $ t \geq 0 $,

\begin{equation}\label{DescentBoundEq}
F\left(x_{t}\right)-F\left(x^{*}\right) \leq \frac{1}{2} L\left[\delta_{\sigma}+\left(1-\frac{\tau}{16 L}\right)^{t}\left(\left\|x_{0}-x^{*}\right\|^{2}-\delta_{\sigma}\right)\right] .
\end{equation}

Here,

\begin{equation}\label{DescentBoundThmThings}
\delta_{\sigma}=\left(\frac{128}{\tau^{2}}+\frac{16}{\tau L}\right) L^{2} d \sigma^{2}+\left(\frac{8}{\tau^{2}}+\frac{1}{2 \tau L}\right) \frac{C^{2}(M !)^{2} \pi d}{4^{M}((2 M) !)^{2}} \sigma^{2} .
\end{equation}

\end{thm}

\begin{proof}

Firstly, let's find an upper bound for $\left\|\widetilde{\nabla}_{\sigma, \Xi}^{M}[F](x)\right\|$. Recall that

\begin{equation}\label{BoundDGSEStimatorSquaredEq}
\left\|\widetilde{\nabla}_{\sigma, \Xi}^{M}[F](x)\right\|^{2} = \sum_{i=1}^d \left|\widetilde{\mathcal{D}}^{M}\left[G_{\sigma}\left(0 \mid x, \xi_{i}\right)\right]\right|^{2}.
\end{equation}

Each term of the sum defined above can be bounded as

\begin{equation}\label{BoundOnErrorEstimationGradientSmoothingSectionEq2}
 \begin{aligned}
 \left|\widetilde{\mathcal{D}}^{M}\left[G_{\sigma}\left(0 \mid x, \xi_{i}\right)\right]\right|^{2} 
 & =  \left|\widetilde{\mathcal{D}}^{M}\left[G_{\sigma}\left(0 \mid x, \xi_{i}\right)\right] - 0 \right|^{2} \\
 & \leq  2\left|\widetilde{\mathcal{D}}^{M}\left[G_{\sigma}\left(0 \mid x, \xi_{i}\right)\right] - \mathcal{D}\left[G_{\sigma}\left(0 \mid x, \xi_{i}\right)\right] \right|^{2} + \\ & \quad 2\left|\widetilde{\mathcal{D}}^{M}\left[G_{\sigma}\left(0 \mid x, \xi_{i}\right)\right] - 0 \right|^{2} \\
 & = 2\left|\widetilde{\mathcal{D}}^{M}\left[G_{\sigma}\left(0 \mid x, \xi_{i}\right)\right] - \mathcal{D}\left[G_{\sigma}\left(0 \mid x, \xi_{i}\right)\right] \right|^{2} + \\ & \quad 2\left|\widetilde{\mathcal{D}}^{M}\left[G_{\sigma}\left(0 \mid x, \xi_{i}\right)\right] \right|^{2} \\
 & \leq \frac{2 C^{2}(M !)^{2} \pi}{4^{M}((2 M) !)^{2}} \sigma^{4 M-2} + 4\left|\widetilde{\mathcal{D}}^{M}\left[G_{\sigma}\left(0 \mid x, \xi_{i}\right)\right] - \nabla_{\xi_{i}} F(x) \right|^{2} + \\ & \quad 4\left|\nabla_{\xi_{i}} F(x)\right|^{2} \\
 & \leq \frac{2 C^{2}(M !)^{2} \pi}{4^{M}((2 M) !)^{2}} \sigma^{4 M-2} + 64\sigma^2 L^2 + 4\left|\nabla_{\xi_{i}} F(x)\right|^{2},
 \end{aligned} 
\end{equation}

where \ref{BoundOnErrorGradientSmoothingSectionEq} and \ref{BoundErrorGHEq} were used. This bound can then used in \ref{BoundDGSEStimatorSquaredEq}, in order to obtain

\begin{equation}\label{BoundDGSEStimatorSquaredEq2}
 \left\|\widetilde{\nabla}_{\sigma, \Xi}^{M}[F](x)\right\|^{2} \leq \frac{2 C^{2}(M !)^{2} d \pi}{4^{M}((2 M) !)^{2}} \sigma^{4 M-2} + 4 \sum_{i=1}^{d}\left|\nabla_{\xi_{i}} F(x)\right|^{2} +  64 d \sigma^2 L^2.
\end{equation}

Now, let's define that $r_t = \norm{x_t - x^*}$. Then

\begin{equation}\label{DifferenceIterationMinimizerEq1}
 \begin{aligned} 
 r_{t+1}^{2}
 & = \norm{x_{t}-\lambda \widetilde{\nabla}_{\sigma, \Xi}^{M} [F]\left(x_{t}\right)-x^{*}}^{2} \\
 & = r_{t}^{2}-2 \lambda\left\langle\widetilde{\nabla}_{\sigma, \Xi}^{M}[F]\left(x_{t}\right), x_{t}-x^{*}\right\rangle+\lambda^{2}\norm{\tilde{\nabla}_{\sigma, \Xi}^{M}[F]\left(x_{t}\right)}^{2} \\ 
 & \stackrel{(\ref{BoundDGSEStimatorSquaredEq2})}{\leq} r_{t}^{2}-2 \lambda\left\langle\widetilde{\nabla}_{\sigma, \Xi}^{M}[F]\left(x_{t}\right)-\nabla F\left(x_{t}\right), x_{t}-x^{*}\right\rangle-2 \lambda\left\langle\nabla F\left(x_{t}\right), x_{t}-x^{*}\right\rangle \\ 
 & +4 \lambda^{2} \sum_{i=1}^{d}\left|\nabla_{\xi_{i}} F\left(x_{t}\right)\right|^{2}+64 \lambda^{2} L^{2} d \sigma^{2}+\frac{2 C^{2} \lambda^{2}(M !)^{2} \pi d}{4^{M}((2 M) !)^{2}} \sigma^{4 M-2}. 
 \end{aligned} 
\end{equation}

Now, the aim is to bound the right side of \ref{DifferenceIterationMinimizerEq1}. Given the strong convexity of $F$, 

\begin{equation}
 -2 \lambda\left\langle\nabla F\left(x_{t}\right), x_{t}-x^{*}\right\rangle \leq 2 \lambda F\left(x^{*}\right)-2 \lambda F\left(x_{t}\right)-\lambda \tau\left\|x^{*}-x_{t}\right\|^{2} . 
\end{equation}

At the same time it holds \textbf{(Here I don't really grasp how/why the first inequality holds true)},

\begin{equation}
 \begin{aligned} 
 & -2 \lambda\left\langle\tilde{\nabla}_{\sigma, \Xi}^{M}[F]\left(x_{t}\right)-\nabla F\left(x_{t}\right), x_{t}-x^{*}\right\rangle \\ 
 & \leq \frac{2 \lambda}{\tau}\norm{\widetilde{\nabla}_{\sigma, \Xi}^{M}[F]\left(x_{t}\right)-\nabla F\left(x_{t}\right)}^{2}+\frac{\lambda \tau}{2}\norm{x_{t}-x^{*}}^{2} \\ 
 & \stackrel{\text { (\ref{BoundErrorEstimatorEq}) }}{\leq} \frac{4 \lambda}{\tau} \cdot \frac{C^{2} \pi d(M !)^{2}}{4^{M}((2 M) !)^{2}} \sigma^{4 M-2}+\frac{64 \lambda}{\tau} L^{2} d \sigma^{2}+\frac{\lambda \tau}{2}\norm{x_{t}-x^{*}}^{2}.
 \end{aligned} 
\end{equation}

Using an estimated for functions with Lipschitz Continuos Gradients from \ref{TeoremaNesterov}, it gives

\begin{equation}\label{BoundSumModulesEq}
 4 \lambda^{2} \sum_{i=1}^{d}\left|\nabla_{\xi_{i}} F\left(x_{t}\right)\right|^{2}=4 \lambda^{2}\left\|\nabla F\left(x_{t}\right)\right\|^{2} \leq 8 \lambda^{2} L\left(F\left(x_{t}\right)-F\left(x^{*}\right)\right) . 
\end{equation}

Now, joining \ref{DifferenceIterationMinimizerEq1} and \ref{BoundSumModulesEq}, it holds

\begin{equation}
 \begin{aligned} 
 r_{t+1}^{2} 
 & \leq r_{t}^{2}-\left(2 \lambda-8 \lambda^{2} L\right)\left(F\left(x_{t}\right)-F\left(x^{*}\right)\right) \\ 
 & +\left(\frac{64 \lambda}{\tau}+64 \lambda^{2}\right) L^{2} d \sigma^{2}+\left(\frac{4 \lambda}{\tau}+2 \lambda^{2}\right) \frac{C^{2}(M !)^{2} \pi d}{4^{M}((2 M) !)^{2}} \sigma^{4 M-2}. 
 \end{aligned} 
\end{equation}

Given the strong convexity of $F$, \textbf{(Here the problem is that the paper from which I took most of this proposition, didn't have an adaptive step-size or smoothing parameter, hence here they could assume a given value for the step size and then the smoothing parameter to be <1. )}

\end{proof}